%-------------------------------------------------------------
\subsection*{Training details}\label{sec:training}
%-------------------------------------------------------------

\textbf{Data split.}
The 86 scenes in \datasetname{} are partitioned into 68
training scenes and 18 held-out test scenes (an 80/20 split),
sampled uniformly across the four driving directions to avoid
directional bias. Each training unit comprises seven
synchronised 29-frame videos at $1280\!\times\!720$, one per
vehicle camera, paired with the R2V control unit and a
structured text prompt (Methods: \emph{R2V control unit,
R2V generation model}).

\textbf{Optimisation.}
We post-train the control pathway with the rectified-flow
velocity-matching
objective~\cite{liu2023rectifiedflow,agarwal2025cosmos},
using a logit-normal time-step distribution and a shift of 5.
Optimisation uses FusedAdamW~\cite{loshchilov2019adamw} with
$\beta\!=\!(0.9,\,0.999)$, weight decay $10^{-3}$, and a peak
learning rate of $1\!\times\!10^{-4}$ that is cosine-annealed
after a 100-step linear warm-up. Gradients are clipped to an
$\ell_2$-norm of $0.1$. The model is trained for 7{,}000
iterations, at which point the rectified-flow loss on the
held-out test scenes has plateaued.

\textbf{Hardware and parallelism.}
Training is distributed over 8 NVIDIA A100 GPUs. The backbone
parameters are sharded across the 8 ranks using fully sharded
data parallelism (FSDP), while the token stream of each
seven-view clip is split across the same 8 ranks along the
view/time axis through context parallelism. Gradients are
accumulated over 4 micro-steps, yielding an effective batch
size of 4 seven-view clips per optimiser step.
Mixed-precision training follows the Wan~2.1
fp32-accumulation strategy to stabilise the rectified-flow
loss under bf16 activation precision.
